% Part: counterfactuals
% Chapter: introduction
% Section: paradoxes-material

\documentclass[../../../include/open-logic-section]{subfiles}

\begin{document}

\olfileid{cnt}{int}{par}

\olsection{实质条件句的悖论}

最早批评用 $!A \lif !B$ 来符号化英语“如果……那么……”语句的人之一是 C.~I. Lewis（刘易斯）。Lewis 批评的是 Whitehead（怀特海）和 Russell（罗素）在《数学原理》中对实质条件句的使用，他们将 $\lif$ 读作“蕴涵”。Lewis 正确地指出，如果 $\lif$ 意指“蕴涵”，那么任何假命题~$p$ 都蕴涵 $p$ 蕴涵~$q$，因为若~$p$ 为假，则 $p \lif (p \lif q)$ 为真；并且任何真命题~$q$ 也蕴涵 $p$ 蕴涵~$q$，因为若 $q$~为真，则 $q \lif (p \lif q)$ 为真。

逻辑学家当然知道，\emph{蕴涵}（implication），即逻辑语义后承（logical entailment），不是联结词，而是公式或语句之间的关系。因此，我们只需避免将 $\lif$ 读作“蕴涵”，就能防止混淆。\footnote{尽管哲学家们试图通过将其读作“仅当”来避免 Lewis 所强调的混淆，但将“$\lif$”读作“蕴涵”的做法在数学家和计算机科学家中仍很普遍。}只要我们避免这种读法，Lewis 所担忧的那种情况根本就不会出现：即使我们把 $p$ 看作代表一个假的英语句子的符号，$p$ 也不会“蕴涵”$q$。要判定是否 $p \Entails q$，我们必须考虑\emph{所有}赋值；即使我们用 $p$ 来符号化一个恰好为假的句子，$p \Entails/ q$ 也依然成立。

但是，像 Lewis 所举的这类“如果……那么……”语句仍然有些奇怪之处：

\begin{quote}
如果月亮是绿奶酪做的，那么 $2+2=4$。
\end{quote}

以及以下推理：

\begin{quote}
月亮不是绿奶酪做的。因此，如果月亮是绿奶酪做的，那么 $2+2=4$。

$2+2 = 4$。因此，如果月亮是绿奶酪做的，那么 $2+2=4$。
\end{quote}

然而，如果“如果……那么……”仅仅是 $\lif$，那么上述句子为真将毫无问题，这些推理的有效性也毫无问题。

另一个例子涉及重言式 $(!A \lif !B) \lor (!B \lif
!A)$。这意味着，如果你从报纸上随机抽取两个直陈句~$S$ 和 $T$，那么“如果 $S$ 那么 $T$，或者如果 $T$ 那么~$S$”这个句子就应该是真的。

\end{document}